\documentclass[11pt]{article}
\usepackage[utf8]{inputenc}
\usepackage[T1]{fontenc}
\usepackage{amsmath, amssymb, amsthm}
\usepackage{mathtools}
\usepackage{geometry}
\geometry{margin=1in}
\usepackage{hyperref}

\newtheorem{theorem}{Theorem}
\newtheorem{proposition}[theorem]{Proposition}
\newtheorem{lemma}[theorem]{Lemma}
\theoremstyle{definition}
\newtheorem{definition}[theorem]{Definition}
\theoremstyle{remark}
\newtheorem{remark}[theorem]{Remark}

\newcommand{\R}{\mathbb{R}}
\newcommand{\N}{\mathbb{N}}
\newcommand{\E}{\mathbb{E}}
\newcommand{\PP}{\mathbb{P}}
\newcommand{\Var}{\mathrm{Var}}
\newcommand{\Cov}{\mathrm{Cov}}
\newcommand{\F}{\mathcal{F}}
\renewcommand{\d}{\mathrm{d}}

\title{Brownian Motion / Wiener Process}
\author{math-foundations: concept 32}
\date{}

\begin{document}
\maketitle

\section{Definitions}

\begin{definition}[Standard Brownian motion / Wiener process]
A stochastic process $\{W_t\}_{t \geq 0}$ on a probability space
$(\Omega, \F, \PP)$ taking values in $\R$ is called a
\emph{standard Brownian motion} (or \emph{Wiener process}) if:
\begin{enumerate}
  \item $W_0 = 0$ almost surely.
  \item For all $0 \leq s < t$, the increment $W_t - W_s$ is Gaussian:
        $W_t - W_s \sim \mathcal{N}(0,\, t-s)$.
  \item \textbf{Independent increments:} for any
        $0 \leq t_0 < t_1 < \cdots < t_n$, the random variables
        $W_{t_1} - W_{t_0},\ W_{t_2} - W_{t_1},\ \dots,\ W_{t_n} - W_{t_{n-1}}$
        are mutually independent.
  \item The sample path $t \mapsto W_t(\omega)$ is continuous in $t$ for
        $\PP$-almost every $\omega \in \Omega$.
\end{enumerate}
\end{definition}

\begin{remark}
Multidimensional Brownian motion $W_t \in \R^d$ consists of $d$
independent copies $W_t = (W_t^{(1)}, \dots, W_t^{(d)})$, each a standard
real Brownian motion.
\end{remark}

\begin{definition}[Natural filtration]
The \emph{natural filtration} of $\{W_t\}$ is the increasing family of
$\sigma$-algebras $\F_t^W := \sigma(W_s : 0 \leq s \leq t)$. We say a
process is adapted to $\{\F_t^W\}$ if it is measurable with respect to
$\F_t^W$ at each time $t$.
\end{definition}

\section{The covariance function}

The defining property of Brownian motion that powers nearly every
computation in stochastic calculus is its covariance kernel:
$$\E[W_s W_t] = \min(s, t).$$
This makes $\{W_t\}$ a centered Gaussian process; in fact this covariance
\emph{together with} Gaussianity and the right continuity modification
characterizes Brownian motion uniquely in distribution.

\begin{theorem}[Covariance of Brownian motion]\label{thm:cov}
Let $\{W_t\}_{t \geq 0}$ be a standard Brownian motion. Then for all
$s, t \geq 0$,
$$\E[W_s W_t] = \min(s, t).$$
\end{theorem}

\begin{proof}
By symmetry of the claim, we may assume without loss of generality that
$0 \leq s \leq t$, so that $\min(s,t) = s$.

\textbf{Step 1: Decompose using independent increments.}
Write
$$W_t = W_s + (W_t - W_s),$$
which is valid by definition. Substituting into the product gives
$$W_s W_t = W_s \big( W_s + (W_t - W_s) \big) = W_s^2 + W_s(W_t - W_s).$$
Taking expectations and using linearity,
\begin{equation}\label{eq:split}
  \E[W_s W_t] = \E[W_s^2] + \E\big[W_s (W_t - W_s)\big].
\end{equation}

\textbf{Step 2: Evaluate the first term.}
By Definition~1 (applied with the increment $W_s - W_0 = W_s$, since
$W_0 = 0$ a.s.), we have $W_s \sim \mathcal{N}(0, s)$. Hence
$$\E[W_s^2] = \Var(W_s) + (\E[W_s])^2 = s + 0 = s.$$

\textbf{Step 3: The cross term vanishes by independence.}
Again by Definition~1, the increments $W_s - W_0 = W_s$ and $W_t - W_s$
are independent (these correspond to disjoint time intervals $[0,s]$ and
$[s,t]$). Independence implies the expectation of the product is the
product of expectations:
$$\E\big[W_s (W_t - W_s)\big] = \E[W_s] \cdot \E[W_t - W_s] = 0 \cdot 0 = 0,$$
since both factors are centered Gaussians with mean $0$.

\textbf{Step 4: Combine.}
Plugging Steps 2 and 3 back into~\eqref{eq:split},
$$\E[W_s W_t] = s + 0 = s = \min(s, t),$$
which proves the theorem for $s \leq t$. The case $t \leq s$ follows by
symmetry of the inner product, completing the proof.
\end{proof}

\begin{remark}[Why this matters]
Theorem~\ref{thm:cov} immediately yields the variance of any linear
functional of $W$. For a deterministic function $f \in L^2([0,T])$, the
Wiener integral $\int_0^T f(s)\, \d W_s$ is Gaussian with mean $0$ and
variance $\int_0^T f(s)^2\, \d s$ — a direct consequence of the kernel
$\min(s,t)$ and the Itô isometry. This is the foundation on which the
entire theory of stochastic integration is built.
\end{remark}

\section{Further properties (stated)}

\begin{proposition}[Scaling / self-similarity]
For any $c > 0$, the rescaled process $\{c^{-1/2} W_{ct}\}_{t \geq 0}$ is
also a standard Brownian motion. Equivalently,
$\{W_{ct}\} \stackrel{d}{=} \{\sqrt{c}\, W_t\}$ as processes.
\end{proposition}

\begin{proposition}[Martingale property]
The process $\{W_t\}$ is a martingale with respect to its natural
filtration: $\E[W_t \mid \F_s^W] = W_s$ for $s \leq t$. Moreover, both
$W_t^2 - t$ and $\exp(\lambda W_t - \tfrac{1}{2}\lambda^2 t)$
($\lambda \in \R$) are also martingales.
\end{proposition}

\begin{proposition}[Path regularity]
Almost surely, the sample path $t \mapsto W_t$ is continuous but nowhere
differentiable, and has unbounded total variation on every nontrivial
interval. Its quadratic variation on $[0,t]$ equals $t$.
\end{proposition}

\end{document}
